\section{Background and Related Work}
\label{sec:background}

\subsection{Hash Table with Separate Chaining}
\label{subsec:hashtable}

A hash table maps keys to array indices via a hash function $h: U \to \{0, 1, \ldots, m-1\}$, where $U$ is the universe of keys and $m$ is the table capacity. When two distinct keys $k_1 \neq k_2$ satisfy $h(k_1) = h(k_2)$, a \emph{collision} occurs. Separate Chaining resolves collisions by maintaining a linked list at each bucket~\citep{cormen2022}.

Under the \emph{simple uniform hashing assumption} (SUHA), the expected chain length equals the load factor $\alpha = n/m$, yielding an expected lookup time of $\Theta(1 + \alpha)$. When $\alpha$ is kept bounded (e.g., $\alpha \leq 0.75$) via dynamic resizing, lookups remain $O(1)$ amortized~\citep{knuth1998}.

\textbf{Memory Cost.} Each element requires storage for:
\begin{itemize}[nosep]
    \item The key string itself ($\approx 40 + \ell$ bytes, where $\ell$ is the string length),
    \item A linked-list node (object header + key reference + next pointer $\approx 32$ bytes),
    \item A bucket reference in the table array ($8$ bytes on a 64-bit JVM).
\end{itemize}
For typical usernames ($\ell \approx 12$), this totals approximately $\mathbf{100}$~\textbf{bytes per element}.

\subsection{Bloom Filter}
\label{subsec:bloomfilter}

A Bloom Filter~\citep{bloom1970} is a probabilistic data structure that represents a set $S$ of $n$ elements using a bit array $B[0 \ldots m{-}1]$ and $k$ independent hash functions $h_1, h_2, \ldots, h_k$. To insert an element $x$, the filter sets bits $B[h_i(x)] \gets 1$ for all $1 \leq i \leq k$. To query membership, it checks whether all $k$ bits are set: if any bit is $0$, the element is \emph{definitely absent}; if all bits are $1$, the element is \emph{probably present} (with a probability of error equal to the FPR).

\textbf{Optimal Parameters} (Mitzenmacher Formula~\citep{mitzenmacher2017}). Given $n$ elements and a target false positive rate $p$, the optimal bit-array size and number of hash functions are:
\begin{align}
    m &= \left\lceil -\frac{n \ln p}{(\ln 2)^2} \right\rceil \label{eq:optimal_m} \\
    k &= \left\lfloor \frac{m}{n} \cdot \ln 2 \right\rceil \label{eq:optimal_k}
\end{align}

For $p = 0.01$ (1\%), these simplify to $m \approx 9.585 \cdot n$ bits and $k = 7$. The resulting memory cost is:
\[
    \text{Memory} = \frac{m}{8} \approx \frac{9.585n}{8} \approx \mathbf{1.2}~\textbf{bytes per element}
\]

\textbf{No False Negatives.} A correctly implemented Bloom Filter guarantees zero false negatives: if $x \in S$, then $\texttt{lookup}(x)$ always returns \texttt{true}.

\subsection{Double Hashing Technique}
\label{subsec:doublehash}

Kirsch and Mitzenmacher~\citep{kirsch2006} proved that only \emph{two} independent hash functions $h_1$ and $h_2$ are needed to simulate $k$ hash functions without increasing the asymptotic FPR:
\begin{equation}
    g_i(x) = \bigl(h_1(x) + i \cdot h_2(x)\bigr) \bmod m, \quad i = 0, 1, \ldots, k{-}1
    \label{eq:doublehash}
\end{equation}
This technique dramatically reduces computational cost---instead of evaluating $k$ independent hash functions, only two hash computations are required per operation.

\subsection{Related Work}

Broder and Mitzenmacher~\citep{broder2004} provide a comprehensive survey of Bloom Filter applications in networking, including distributed caching, routing, and peer-to-peer systems. Fan et al.~\citep{fan2014} propose the Cuckoo Filter as an alternative that supports deletion, though with higher per-element overhead. Chang et al.~\citep{chang2008} describe the use of Bloom Filters in Google's Bigtable to reduce unnecessary disk reads---a direct application of the crossover phenomenon we study.

Our work differs from prior studies by focusing specifically on the \emph{transition boundary} under strict, quantified memory constraints, using the JVM's \texttt{-Xmx} flag to simulate hard RAM budgets and observe the exact crossover dynamics.
